% Override DesignTestOptions to test another class size/aspect ratio/notheorems.
\providecommand{\DesignTestOptions}{10pt,aspectratio=169}
\documentclass[\DesignTestOptions,xcolor={svgnames}]{beamer}
\usepackage{clean_style_beamer}
\addbibresource{beamer_design_test.bib}
\pdfstringdefDisableCommands{\def\translate#1{#1}}

\title[Design proof]{A Shared Design for Research and Evidence}
\subtitle[Shared style]{Typography, emphasis, and reusable elements}
\author[A. Researcher et al.]{Alexandra Researcher\inst{1} \and
  Benjamin Collaborator\inst{2} \and Cecilia Investigator\inst{1} \and
  Daniel Scholar\inst{3}}
\institute[Research institutes]{\inst{1} Research Institute \quad
  \inst{2} University of Social Sciences \quad \inst{3} Policy Centre}
\date[25 Oct 2026]{25 October 2026}

\makeatletter
\newcommand{\AssertListSpacing}{%
  \begingroup
    \dimen@=\itemsep
    \csname @list\romannumeral\@listdepth\endcsname
    \addtolength{\itemsep}{0.8\baselineskip}%
    \ifdim\dimen@=\itemsep\else
      \PackageError{beamer_design_test}{Incorrect list spacing}{Spacing must be added once at the current list font size.}%
    \fi
  \endgroup}
\@for\TestEnvironment:=theorem,corollary,fact,lemma,problem,solution,definition,definitions,example,examples,exercise,assumption,proposition,Beispiel,Beispiele,Loesung,Satz,Folgerung,Fakt,Beweis,Lemma,Proof,Theorem,Problem,Corollary,Example,Examples,Definition,mblock\do{%
  \@ifundefined{\TestEnvironment}{%
    \PackageError{beamer_design_test}{Missing environment \TestEnvironment}{All public environments must remain available.}}{}}
\makeatother

\begin{document}
\maketitle

\section{Typography and Emphasis}

\begin{frame}{Typography and Semantic Sizes}
  Ordinary body text uses the current Source Serif family.
  Proportional old-style numerals: 0123456789.
  \begin{itemize}
    \item Regular, \textbf{bold}, \textit{italic}, and \textbf{\textit{bold italic}}.
    \item \tbf{Blue bold emphasis} and \alert{coral alert text} are distinct.
      \begin{itemize}
        \item Nested text follows the class-relative size hierarchy.
      \end{itemize}
  \end{itemize}
  \smallcite{\textcite{smith2020}}\quad
  \hyperlink{buttons}{\beamergotobutton{Buttons and links}}
\end{frame}

\begin{frame}{Nested Lists and Overlays}
  \begin{itemize}[<+->]
    \AssertListSpacing
    \item First-level dot-circle; text and marker use relative list layout.
      \begin{itemize}[<.->]
        \AssertListSpacing
        \item Second-level solid long arrow in a fixed-width icon box.
          \begin{itemize}[<.->]
            \AssertListSpacing
            \item The existing third-level marker is preserved.
          \end{itemize}
        \AssertListSpacing
      \end{itemize}
    \AssertListSpacing
    \item Spacing and grouping survive the nested environment.
  \end{itemize}
\end{frame}

\begin{frame}{Numbered Lists and Optional Labels}
  \begin{enumerate}[<+->][a)]
    \AssertListSpacing
    \item Custom alphabetic labels remain blue.
      \begin{enumerate}[<.->][i)]
        \AssertListSpacing
        \item Nested numbering and overlay defaults are accepted.
        \item The extra spacing is applied once.
      \end{enumerate}
    \AssertListSpacing
    \item The outer counter continues after the nested list.
  \end{enumerate}
\end{frame}

\subsection{Reusable Callouts}

\begin{frame}{Shared Emphasis and Callout Helpers}
  \begin{mblock}
    \alert{\textbf{Research question:}} Can one shared preamble preserve a
    clear visual hierarchy across presentations?
  \end{mblock}
  \begin{itemize}
    \item \tbf{Reusable emphasis:} blue, bold serif text.
    \item \alert{Alerts:} coral, with surrounding text staying black.
    \item The callout retains its grey fill and rounded orange border.
  \end{itemize}
\end{frame}

\section*{Starred Section Without an Automatic Divider}
\begin{frame}{Starred Section Behaviour}
  This frame follows a starred section directly. No automatic divider is
  inserted for the starred section.
\end{frame}

\section{Theorems and Equations}

\begin{frame}{A Numbered Proposition with an Optional Name}
  \begin{proposition}<1->[An illustrative identity]
    For every real number $x$, the following identity holds:
    \begin{equation}\label{eq:identity}
      (x+1)^2=x^2+2x+1.
    \end{equation}
    Plain-style theorem prose is italic; the heading is blue and numbered.
  \end{proposition}
  \begin{proof}
    Expand the product $(x+1)(x+1)$.
  \end{proof}
\end{frame}

\begin{frame}{Theorem Counters and Compatibility Aliases}
  \begin{theorem}[First named result]
    The native theorem counter remains available.
  \end{theorem}
  \begin{Lemma}[A compatible alias]
    Existing aliases retain optional names and their underlying counters.
  \end{Lemma}
  \begin{proposition}[Second proposition]
    Proposition numbering advances automatically.
  \end{proposition}
\end{frame}

\begin{frame}{Definition and Example Semantics}
  \begin{definition}[A reusable design]
    A shared style supplies presentation-wide formatting.
  \end{definition}
  \begin{example}[Upright example prose]
    Definitions and examples retain their normal theorem-style semantics.
  \end{example}
  \begin{assumption}
    Assumptions retain plain-style italic prose.
  \end{assumption}
\end{frame}

\section{References and Captions}

\begin{frame}{Citation Names, Colours, and Local Grouping}
  Ordinary author--year links: \textcite{smith2020,smith2021}.
  \begin{itemize}
    \item Reduced grey citation: \smallcite{\textcite{smith2020}}.
    \item Ordinary links return to magenta: \citep{smith2021}.
    \item Three authors are shortened: \citet{adams2022}.
    \item Author and year commands: \citeauthor{smith2020},
      \citeyearpar{smith2020}.
  \end{itemize}
  Surnames stay compact even when authors share a surname.
\end{frame}

\begin{frame}{Numbered Figures and Subcaptions}
  \begin{figure}
    \centering
    \begin{subfigure}{0.45\linewidth}
      \centering\fbox{\strut First panel}
      \caption{First descriptive panel caption.}
    \end{subfigure}\hfill
    \begin{subfigure}{0.45\linewidth}
      \centering\fbox{\strut Second panel}
      \caption{Second descriptive panel caption.}
    \end{subfigure}
    \caption{A deliberately long figure caption that wraps naturally within
      its containing width while keeping its numbered label blue and its
      descriptive text black.}\label{fig:panels}
  \end{figure}
\end{frame}

\begin{frame}{Booktabs Table with a Numbered Caption}
  \begin{table}
    \caption{An illustrative table with automatic numbering.}\label{tab:sample}
    \centering
    \begin{tabular}{lrr}
      \toprule
      Group & Baseline & Alternative \\
      \midrule
      Primary & 12.3 & 14.5 \\
      Secondary & 23.4 & 25.6 \\
      Tertiary & 34.5 & 36.7 \\
      \bottomrule
    \end{tabular}
  \end{table}
  Table text retains the ordinary class-relative body size.
\end{frame}

\begin{frame}[label=buttons]{Buttons and Hyperlinks}
  \hyperlink{appendix-proof}{\beamergotobutton{Appendix details}}\quad
  \hyperlink{appendix-proof}{\beamerreturnbutton{Return-style button}}
  \begin{itemize}
    \item Buttons have white fill, blue text/borders, and tiny monospace labels.
    \item External URL: \href{https://example.org}{example.org}.
    \item Local reference: equation~\eqref{eq:identity}.
    \item Figure~\ref{fig:panels} and table~\ref{tab:sample} use automatic counters.
  \end{itemize}
\end{frame}

\begin{frame}[fragile]{Existing Code Listing Support}
\begin{lstlisting}
% An illustrative Matlab calculation
x = ones(2, 1);
y = factorial(3) * x;
\end{lstlisting}
  Listing colours and monospace behaviour retain their existing settings.
\end{frame}

\begin{frame}[allowframebreaks]{Bibliography}
  \printbibliography[heading=none]
\end{frame}

\appendix
\section{Appendix}
\begin{frame}[label=appendix-proof]{Appendix and Return Navigation}
  Appendix-numbering support remains active with no visible slide numbers.
  \begin{equation}
    \mathbb{E}[X]=\int_{\mathbb{R}} x f(x)\,\mathrm{d}x.
  \end{equation}
  \hyperlink{buttons}{\beamerreturnbutton{Return to buttons}}
\end{frame}
\end{document}
